\section{Quantization speculation and structured action outputs}\label{sec:C31:05}
Weight quantization, activation quantization, and KV quantization change different error pathways. GPTQ targets post-training weight approximation, while SmoothQuant redistributes activation difficulty to make low-precision execution more practical~\cite{frantar2022gptq,xiao2022smoothquant}. Their historical methods remain useful foundations; the current implementation choice must be checked against the target architecture, kernels, and 2026 model release. A nominal bit width is not a complete specification: group size, scales, accumulation precision, outlier handling, and supported matrix shapes affect both numerical behavior and performance.

For an agent, a small logit perturbation can cause a discrete tool change. Average text perplexity is therefore an incomplete acceptance metric. Evaluate schema validity, tool selection, argument fidelity, termination, rare constraints, and task success. If the top two action probabilities are close, a modest numerical change can reverse their ordering even when most token probabilities are nearly unchanged. A quality-preserving quantization claim needs the decision distribution and workload that matter to the application.

Speculative decoding proposes tokens with a draft mechanism and verifies them with the target model. Under the appropriate acceptance rule, the goal is to preserve the target sampling distribution while reducing costly sequential target work~\cite{leviathan2022speculative}. Its benefit depends on acceptance, verification cost, draft cost, batch interaction, and available hardware parallelism. A high acceptance rate does not guarantee a speedup if the draft consumes the same bottleneck resource or increases memory pressure enough to reduce concurrency.

Let a draft propose $k$ tokens, with expected accepted progress $G_k$, draft time $D_k$, and target verification time $V_k$. Ignoring overlap, the useful condition relative to target-only per-token time $t$ is
\begin{equation}
 \frac{D_k+V_k}{G_k}<t.
\end{equation}
For independent identically distributed cycles with finite mean duration and positive finite mean progress, the ratio of those means is the long-run time per emitted token by renewal-reward accounting. It is neither the expectation of each cycle's time/progress ratio nor a latency percentile. Adaptive cycles, shared-resource contention, and changing context lengths require a justified stationary model or direct measurement. Structured tool outputs create additional boundaries. Drafted tokens cannot trigger side effects before the target output, parser, authorization policy, and commit logic accept the action. Streaming a plausible partial tool call into execution defeats the intended separation between speculative text and committed behavior.

